\section{从电路到操作系统：组成原理的最小主线}

很多读者第一次学操作系统时，会被一组硬件词汇挡住：CPU、寄存器、总线、MMU、DMA、APIC、IOMMU、cache、TLB、特权级。若只把这些词当作背景知识，后面的系统调用、虚拟内存、调度和驱动都会像悬空的规则。本节把计算机组成原理和操作系统接起来：先从最小数字电路开始，说明 CPU 和设备为什么这样设计，再解释 OS 为什么必须依赖这些硬件机制。

\begin{keyidea}
操作系统不是硬件之外的魔法层。它是运行在受保护 CPU 模式下的一段程序，利用中断、异常、页表、定时器、总线、设备寄存器和 DMA，把原始硬件状态组织成进程、地址空间、文件和 socket。
\end{keyidea}

\subsection{从开关到位：硬件最初解决什么问题}

最底层的电子计算机不是从“进程”开始的，而是从可控电信号开始的。晶体管可以近似看作一个由输入电压控制的开关。开关组合起来可以实现逻辑门，逻辑门组合起来可以实现加法器、比较器、选择器和译码器。把电压高低抽象成 1 和 0，就得到数字电路的基本单位：bit。

\begin{longtable}{p{0.18\linewidth}p{0.34\linewidth}p{0.36\linewidth}}
\toprule
层次 & 解决的问题 & 向上一层暴露的抽象 \\
\midrule
晶体管 & 电信号是否导通 & 可控开关 \\
逻辑门 & 布尔函数如何计算 & AND、OR、NOT、XOR \\
组合逻辑 & 输入如何变成输出 & 加法器、比较器、译码器、MUX \\
触发器 & 如何记住一个 bit & 时钟边沿采样的状态 \\
寄存器 & 如何保存一组 bit & 机器字、地址、临时值 \\
有限状态机 & 如何按步骤推进 & 当前状态、输入、下一状态 \\
\bottomrule
\end{longtable}

组合逻辑没有记忆。给定输入，它经过门延迟后产生输出；输入一变，输出也会变。若只靠组合逻辑，机器不能记住“上一步执行到哪里”。触发器和寄存器引入记忆：在时钟边沿，把输入采样到内部状态，直到下一次时钟边沿才改变。时钟让复杂电路分成离散步骤，状态机让机器有“现在处于哪一步”的概念。

一个最小计数器已经包含操作系统后来要关心的三个事实：状态存在哪里，什么时候改变，谁能改变。计数器的状态在寄存器里；时钟边沿触发改变；使能信号决定是否更新。把这个思想放大，就得到 CPU 的程序计数器、寄存器文件、状态寄存器和控制逻辑。

\begin{lstlisting}
on each clock edge:
  if reset:
    counter = 0
  else if enable:
    counter = counter + 1
\end{lstlisting}

这段伪代码看起来像软件，实际对应硬件状态机。操作系统学习中反复出现的“状态机”不是写作技巧，而是硬件和内核共同的基本形态：硬件在时钟和事件下推进，内核对象在系统调用、中断和锁保护下推进。

\subsection{从固定电路到存储程序}

早期最简单的计算电路可以为某个特定任务连线，比如把两个输入相加，或者按固定流程控制设备。但一台通用计算机不能每换一个程序就重新接线。存储程序思想把“要做什么”也编码成内存里的数据：指令存放在内存中，CPU 取出指令、解释指令、执行指令，再取下一条。

\begin{lstlisting}
while true:
  inst = memory[rip]
  rip = rip + instruction_length
  decode inst
  execute inst
\end{lstlisting}

这就是 CPU 最小执行循环。它需要至少四类状态：

\begin{itemize}
\tightlist
\item \textbf{程序计数器}：下一条指令在哪里。
\item \textbf{寄存器文件}：当前计算的临时值、地址和参数。
\item \textbf{状态寄存器}：条件码、中断使能、当前模式等控制位。
\item \textbf{内存接口}：从地址读取指令和数据，把结果写回内存。
\end{itemize}

存储程序模型引入了一个强大的统一视角：程序是内存中的字节，数据也是内存中的字节，CPU 根据指令格式解释这些字节。后面的 ELF 装载、地址空间、页表权限、代码段只读可执行、数据段可写不可执行，都是在管理这组字节的含义和权限。

\subsection{ISA 和微架构：程序员看到的机器与芯片内部的机器}

学习 CPU 时必须区分 ISA 和微架构。ISA 是指令集架构，是软件可见的契约：有哪些寄存器，指令如何编码，内存访问怎么表达，异常如何报告，系统调用如何进入内核。微架构是芯片内部怎样实现这个契约：流水线几级，是否乱序执行，cache 多大，分支预测怎么做，指令如何拆成微操作。

\begin{longtable}{p{0.2\linewidth}p{0.34\linewidth}p{0.34\linewidth}}
\toprule
概念 & 例子 & OS 关心方式 \\
\midrule
ISA & x86-64 & 系统调用 ABI、异常入口、页表格式、Ring 0/3 特权级 \\
微架构 & 流水线、乱序、cache、预测器 & 性能、屏障、侧信道、上下文切换成本 \\
ABI & 参数寄存器、栈对齐、返回值 & 用户态和内核态、程序和库之间的二进制契约 \\
\bottomrule
\end{longtable}

同一个 x86-64 程序可以在不同厂商、不同代际 CPU 上运行，是因为这些 CPU 都承诺实现同一套软件可见 ISA。它们内部可能完全不同：有的能同时发射更多微操作，有的分支预测更强，有的 cache 更大。OS 既要遵守 ISA 契约，又要理解微架构成本。例如页表格式来自 ISA，TLB shootdown 成本和 cache 抖动则来自微架构和多核实现。

\subsection{数据通路：一条整数加法如何变成硬件动作}

把 CPU 想象成“会执行代码的盒子”太粗。更有解释力的方式是看数据通路。以 Intel 语法的 \asmcode{add rax, rbx} 为例，CPU 需要从寄存器文件读 \code{rax/rbx}，送入 ALU，加法结果写回 \code{rax}。控制逻辑根据 opcode 产生读寄存器编号、ALU 操作码、写回使能等信号。

\begin{lstlisting}
decode add rax, rbx:
  a = regfile.read(rax)
  b = regfile.read(rbx)
  y = alu.add(a, b)
  regfile.write(rax, y)
\end{lstlisting}

再看 \asmcode{mov rax, qword ptr [rbx + imm]}。它不仅用地址生成单元计算地址，还要访问内存层次，并可能触发缺页异常。若地址合法且 cache 命中，结果很快写回寄存器；若 TLB 未命中，硬件会执行页表遍历；若页不存在或权限不对，CPU 不能继续执行这条指令，必须进入异常入口。

\begin{lstlisting}
decode mov rax, qword ptr [rbx + imm]:
  va = regfile.read(rbx) + imm
  pa = translate(va, READ, current_privilege)
  if translation_fault:
    raise page_fault
  data = cache_or_memory_read(pa)
  regfile.write(rax, data)
\end{lstlisting}

这说明缺页异常不是 OS 主动“检查每次内存访问”。硬件在普通指令执行途中发现不能完成架构语义，于是转入内核。虚拟内存的速度和安全性都来自这个硬件检查路径。

\subsection{为什么需要流水线}

如果 CPU 完整执行完一条指令才开始下一条，许多硬件部件会空闲。流水线把取指、译码、执行、访存、写回分成阶段，不同指令同时处在不同阶段，就像工厂流水线上不同工位同时处理不同产品。

\begin{lstlisting}
cycle 1: inst1 fetch
cycle 2: inst1 decode   inst2 fetch
cycle 3: inst1 execute  inst2 decode   inst3 fetch
cycle 4: inst1 memory   inst2 execute  inst3 decode
cycle 5: inst1 writeback inst2 memory  inst3 execute
\end{lstlisting}

流水线带来吞吐，也带来 hazard。数据 hazard 是后一条指令需要前一条尚未写回的结果；控制 hazard 是分支决定下一条 RIP；结构 hazard 是两个阶段争同一硬件资源。硬件用转发、暂停、预测和乱序执行缓解这些问题。

OS 为什么关心流水线？第一，中断和异常必须是精确的：CPU 要能报告“哪条指令出错，之前的指令已经提交，之后的指令没有对架构状态产生影响”。没有精确异常，内核无法可靠处理 page fault 和信号。第二，分支预测、乱序执行和缓存会留下微架构痕迹，安全隔离必须考虑侧信道。第三，内存屏障的意义来自 CPU 可能为了性能重排某些可见顺序。

\subsection{从单核到多核：并行不是免费加速}

单核 CPU 同一时刻只执行一个控制流。多核把多个 CPU 核放在同一系统里，让多个线程可以真正同时运行。多核增加吞吐，也把“共享内存”变成硬件协议问题：同一个物理地址的 cache line 可能同时存在于多个核心的 cache 中，硬件一致性协议要保证普通读写能以某种规则传播。

\begin{longtable}{p{0.2\linewidth}p{0.34\linewidth}p{0.34\linewidth}}
\toprule
机制 & 解决的问题 & OS 设计影响 \\
\midrule
cache coherence & 多核缓存同一地址如何保持一致 & 锁竞争、伪共享、per-CPU 数据 \\
atomic RMW & 读改写不能被其他核插入 & 自旋锁、引用计数、无锁队列 \\
IPI & CPU 之间发送控制事件 & 重调度、TLB shootdown、停止其他 CPU \\
NUMA & 内存到不同 CPU 距离不同 & CPU affinity、first-touch、页迁移 \\
\bottomrule
\end{longtable}

多核系统里，关本核中断只能阻止本核被中断处理函数打断，不能阻止另一个核访问同一数据。因此从单片机临界区演进到 SMP 内核时，必须引入自旋锁、原子操作、内存屏障和 per-CPU 数据结构。调度器也不能只维护一个全局队列，因为全局队列会变成锁和 cache line 热点。

\subsection{内存芯片、地址线和最早的内存模型}

CPU 要读写内存，最朴素的硬件接口是地址线、数据线和控制线。CPU 把地址放到地址线上，拉高读信号，内存芯片把对应数据放到数据线上；写入时 CPU 把地址和数据放好，再拉高写信号。真实现代内存控制器复杂得多，但这个模型能解释地址、总线和 MMIO 的根。

\begin{lstlisting}
memory_read(address):
  address_bus = address
  control.read = 1
  wait until memory_ready
  return data_bus

memory_write(address, value):
  address_bus = address
  data_bus = value
  control.write = 1
  wait until memory_ready
\end{lstlisting}

若某个地址范围接的是 RAM，读写就是普通内存。若某个地址范围接的是设备寄存器，读写就变成控制硬件。这就是 memory-mapped I/O 的根本含义：CPU 使用同一套 load/store 机制访问内存和设备，但设备寄存器有副作用，不能当普通内存优化。

\subsection{总线为什么出现：多个部件如何共享通信通道}

当系统里只有 CPU 和一片内存时，连线可以很直接。加入定时器、串口、磁盘控制器、网卡、显示控制器后，若每个部件都和 CPU 单独连一套线，系统会失控。总线把地址、数据和控制协议标准化，让多个主设备和从设备共享通信通道。

\begin{longtable}{p{0.22\linewidth}p{0.34\linewidth}p{0.32\linewidth}}
\toprule
总线概念 & 含义 & OS 可见影响 \\
\midrule
地址译码 & 某个地址范围属于哪个设备 & ACPI、PCI BAR、MMIO 映射 \\
仲裁 & 多个发起者谁先使用总线 & DMA 与 CPU 竞争带宽 \\
事务 & 一次读写请求和响应 & posted write、读回确认、错误状态 \\
中断消息 & 设备如何通知 CPU & IRQ、MSI/MSI-X、中断亲和性 \\
\bottomrule
\end{longtable}

早期系统常有前端总线、南北桥、独立 I/O 总线；现代系统把内存控制器、PCIe root complex、片上互联和 I/O MMU 集成得更深。拓扑变了，OS 面对的问题没变：发现设备、映射寄存器、分配中断、建立 DMA 地址空间、处理错误和热插拔。

\subsection{MMIO 和普通内存的根本区别}

MMIO 地址看起来也是地址，但读写语义不同。读一个状态寄存器可能清除某个 pending 位；写一个 doorbell 寄存器可能让设备开始读取描述符；连续两次写寄存器的顺序可能是设备协议的一部分。编译器不能删除“看似没用”的 MMIO 读，CPU 不能随意把 doorbell 写提前到描述符写之前。

\begin{lstlisting}
// 正确顺序：先准备内存中的命令，再通知设备
desc.addr = dma_address
desc.len = length
desc.flags = READ
dma_wmb()
writel(queue_tail, dev.DOORBELL)
\end{lstlisting}

这里的 \code{dma\_wmb} 不是装饰。它表达“设备看到 doorbell 之前，必须能看到描述符内容”。如果缺少屏障，CPU 或互联可能让设备先收到通知，再读到半初始化描述符。驱动 bug 往往不是每次复现，因为顺序竞态受 CPU、互联、cache 和负载影响。

\subsection{DMA 为什么出现：CPU 不应该搬每一个字节}

PIO 模型里，CPU 亲自从设备寄存器读字节再写到内存，或者反过来。少量控制数据可以这么做，大量磁盘块、网络包、音视频数据不能这么做，否则 CPU 大量时间浪费在搬运上。DMA 让设备成为总线主设备，直接读写内存；CPU 只负责准备 buffer、描述符和地址映射。

\begin{lstlisting}
pio_read_block():
  for i in 0..block_size-1:
    memory[buf+i] = read_device_data_register()

dma_read_block():
  desc.addr = dma_map(buf)
  desc.len = block_size
  desc.cmd = READ
  dma_wmb()
  writel(desc_index, dev.DOORBELL)
  sleep until completion interrupt
\end{lstlisting}

DMA 的性能来自异步和批量，复杂性也来自异步和批量。提交 DMA 后，buffer 的所有权暂时交给设备；completion 前不能释放、复用或移动。设备写入内存后，CPU cache 可能需要失效才能看到新数据。若系统有 IOMMU，内核还要为设备建立 DMA 地址到物理页的映射；若没有 IOMMU，设备 bug 或恶意设备可能写坏任意内存。

\begin{longtable}{p{0.22\linewidth}p{0.34\linewidth}p{0.32\linewidth}}
\toprule
DMA 阶段 & 内核必须做什么 & 常见失败 \\
\midrule
准备 buffer & 分配或 pin 页，确认方向和长度 & 用户页被换出或提前释放 \\
建立映射 & \code{dma\_map} 或 IOMMU map & 地址超出设备 DMA mask \\
发布描述符 & 写描述符并执行 DMA 屏障 & 设备读到旧描述符 \\
等待完成 & 中断、轮询或 completion 队列 & 丢中断、超时、中断风暴 \\
回收资源 & unmap、释放 buffer、唤醒等待者 & use-after-free、重复 completion \\
\bottomrule
\end{longtable}

理解 DMA 后，块层、网络栈和 io\_uring 的很多设计就不再神秘。它们都在围绕同一件事组织队列：把 CPU 上的请求变成设备能批量消费的描述符，把设备完成再变成内核和用户态能理解的 completion。

\subsection{为什么需要中断控制器}

设备事件不能全靠 CPU 轮询，否则空闲时浪费 CPU，高负载时又容易响应不及时。中断让设备在事件发生时通知 CPU。但设备很多、CPU 也可能很多，系统需要一个控制层决定：哪个设备的事件送到哪个 vector，哪个 CPU 接收，优先级如何，中断是否屏蔽，处理完成后如何确认。

\begin{lstlisting}
device raises interrupt
  -> interrupt controller records pending event
  -> controller selects target CPU and vector
  -> CPU finishes current instruction boundary
  -> CPU enters interrupt gate
  -> kernel handler reads device status
  -> kernel sends EOI
\end{lstlisting}

在 x86-64 上，本地 APIC、IOAPIC、x2APIC、MSI/MSI-X 和 IPI 共同完成这些工作。OS 关心的问题包括中断投递、屏蔽、优先级、亲和性、结束通知和错误统计。

\subsection{特权级不是口号：它是 CPU 执行规则的一部分}

“特权”不是说内核更重要，而是 CPU 在硬件上维护的一组模式位和访问规则。当前模式决定哪些指令能执行、哪些寄存器能改、哪些地址能访问、哪些异常入口能跳转。用户态不能随便关中断、改页表、写控制寄存器、访问设备 MMIO，也不能跳到任意内核地址。

\begin{longtable}{p{0.22\linewidth}p{0.32\linewidth}p{0.34\linewidth}}
\toprule
层级 & x86-64 角色 & 说明 \\
\midrule
Ring 3 & 用户程序 & 不能执行特权指令，不能访问 supervisor-only 页和内核 MMIO \\
Ring 0 & 内核 & 可以配置页表、中断、MSR、设备资源和调度状态 \\
VMX root & hypervisor & 可运行和控制 guest，拦截敏感事件并管理二级地址翻译 \\
\bottomrule
\end{longtable}

特权级解决的是“不可信程序不能直接控制整台机器”。没有它，一个普通程序可以改页表读取别的进程，可以关中断让调度器失效，可以写磁盘控制器覆盖文件系统，可以改定时器破坏时间。OS 的安全边界首先依赖 CPU 拒绝这些操作。

\subsection{从没有保护到有保护：一步一步推出来}

最小单片机通常没有用户态/内核态之分。所有代码同权，任何函数都能访问所有寄存器和内存。这在单一固件里可接受，因为所有代码由同一开发者构建，故障域就是整台设备。一旦机器要运行多个程序，问题立即出现。

\begin{enumerate}
\tightlist
\item 多个程序共享内存：程序 A 写坏程序 B 的数据，系统无法定位和阻止。
\item 程序可能死循环：若没有定时器抢占，其他程序永远得不到 CPU。
\item 程序直接碰设备：一个程序可以重置网卡、覆盖磁盘、窃取键盘输入。
\item 程序能改控制寄存器：它可以关闭隔离机制本身。
\end{enumerate}

于是硬件逐步增加四类保护能力。第一是特权级，限制敏感指令。第二是异常入口，让非法操作转入内核处理，而不是让机器随机继续。第三是 MMU 和页表，为每个程序定义可访问地址范围和权限。第四是定时器中断，让内核即使在程序不合作时也能收回 CPU。

\begin{longtable}{p{0.2\linewidth}p{0.34\linewidth}p{0.34\linewidth}}
\toprule
失败 & 硬件机制 & OS 抽象 \\
\midrule
程序互相写内存 & MMU、页表、权限位 & 进程地址空间 \\
程序执行特权操作 & privilege mode、illegal instruction trap & 系统调用和权限检查 \\
程序长期占 CPU & timer interrupt & 抢占式调度 \\
程序直接访问设备 & MMIO 只映射给内核、IOMMU & 驱动、fd、设备文件 \\
\bottomrule
\end{longtable}

这条演进线很重要。特权级、页表、中断和系统调用不是先背定义再硬套，它们分别修复了无保护多程序系统中的具体失败。

\subsection{系统调用为什么不能只是函数调用}

函数调用只是在当前特权级内跳到另一个地址。若用户程序能直接调用内核内部函数，就必须允许它跳进内核地址、使用内核栈、传入任意指针并修改内核状态。这等于没有隔离。系统调用的本质是受控陷入：用户程序只能从固定入口进入，CPU 自动切换到更高特权，内核保存用户上下文，检查参数和权限，再返回用户态。

\begin{lstlisting}
user code:
  put syscall number and args in ABI registers
  execute syscall/trap instruction

hardware:
  verify transition target
  switch privilege mode
  save return RIP and RFLAGS
  jump to kernel entry

kernel:
  switch to kernel stack
  build trap frame
  validate args and permissions
  run syscall implementation
  restore user state and return
\end{lstlisting}

这解释了为什么系统调用有成本，也解释了为什么 vDSO、io\_uring、批量系统调用和共享内存队列能优化某些路径：它们不是绕开安全，而是在不破坏边界的前提下减少陷入次数或把安全检查前移。

\subsection{MMU 的演进：从基址界限到多级页表}

最简单的内存保护可以是基址和界限寄存器：程序访问地址时，硬件检查地址是否落在 \code{base <= addr < base + limit}。这能阻止越界，但无法高效表达稀疏地址空间、共享库、按页换入换出、不同页面不同权限。页表把地址空间切成固定大小的页，每页有独立映射和权限。

\begin{lstlisting}
translate(virtual_address, access):
  vpn = virtual_address / PAGE_SIZE
  off = virtual_address % PAGE_SIZE
  pte = page_table[vpn]
  if !pte.present:
    raise page_fault
  if access not allowed by pte.permissions:
    raise protection_fault
  return pte.physical_page * PAGE_SIZE + off
\end{lstlisting}

多级页表是空间优化。64 位虚拟地址空间很大，若为每个虚拟页都准备一项线性页表，绝大多数项会浪费。多级页表只为实际使用的地址区间分配下级表。TLB 则是时间优化：页表遍历太慢，所以 CPU 缓存最近的地址翻译。

\begin{longtable}{p{0.22\linewidth}p{0.34\linewidth}p{0.32\linewidth}}
\toprule
页表能力 & 直接效果 & OS 用途 \\
\midrule
present 位 & 页面是否有物理页 & 懒分配、文件映射、换页 \\
R/W/X 权限 & 读写执行控制 & 代码只读、栈不可执行、COW \\
user/supervisor & 用户态是否可访问 & 用户/内核地址隔离 \\
dirty/accessed & 是否写过或访问过 & 回收、writeback、老化算法 \\
global/PCID/ASID & TLB 保留和地址空间标记 & 降低上下文切换成本 \\
\bottomrule
\end{longtable}

页表不是“虚拟地址到物理地址的 map”这么简单。它是隔离、懒加载、共享、权限、安全和性能共同使用的硬件数据结构。

\subsection{cache 和 TLB：为什么抽象有成本}

内存比 CPU 慢得多，页表遍历也比寄存器操作慢得多。cache 缓存数据和指令，TLB 缓存地址翻译。它们让程序跑得快，也让 OS 付出维护成本。切换地址空间时，旧 TLB 项不能错误用于新进程；修改页表后，其他 CPU 的旧翻译可能仍在 TLB 中；共享数据频繁写入时，cache line 会在 CPU 间来回迁移。

\begin{lstlisting}
unmap_page(process, va):
  remove PTE from page table
  for each cpu running this address space:
    send IPI to invalidate TLB entry
  wait for acknowledgements
  only now free physical page
\end{lstlisting}

如果跳过 TLB shootdown，某个 CPU 可能继续用旧翻译访问已经释放的物理页。若这页后来分配给内核对象或另一个进程，隔离就被破坏。很多 OS 机制看起来“为什么这么麻烦”，答案往往在 cache/TLB 和多核一致性里。

\subsection{精确异常：OS 能恢复执行的前提}

当一条指令触发缺页或非法操作时，内核需要知道它是否可以修复并返回继续执行。例如用户访问尚未分配但合法的栈页，内核可以分配物理页、更新 PTE，然后返回让同一条 load/store 重试。要做到这一点，CPU 必须提供精确异常：异常点之前的指令已经按顺序提交，异常点之后的指令没有对架构状态产生影响。

\begin{lstlisting}
user instruction stream:
  inst1 committed
  inst2 committed
  inst3 causes page fault  <-- saved RIP points here
  inst4 not architecturally committed

kernel:
  inspect fault address and error code
  fix mapping or deliver signal
  return to inst3 if fixed
\end{lstlisting}

乱序 CPU 内部可能已经执行了 inst4 的某些微操作，但在架构层必须让 OS 看到一个干净边界。没有这个契约，page fault、debug exception、signal、single step 和进程恢复都难以实现。

\subsection{从芯片到主板：固件为什么存在}

真实机器不是只有 CPU 和内存。还有 DRAM 初始化、时钟、PCIe 拓扑、ACPI 表、IOMMU、显卡、存储控制器、键盘、TPM、电源管理和多核启动。CPU 上电时，很多硬件尚未处在 OS 可直接使用的状态。固件负责完成最早期平台初始化，并把硬件描述交给 bootloader 或内核。

\begin{lstlisting}
power on
  -> reset vector
  -> firmware initializes DRAM and chipset
  -> firmware discovers platform tables
  -> bootloader or EFI stub loads kernel
  -> kernel receives memory map, initrd, command line, ACPI tables
  -> kernel takes over devices and memory management
\end{lstlisting}

固件不是万能可信源。内核仍要校验内存地图、保留区、ACPI 描述和启动参数。若固件把内核占用页标为空闲，页分配器会覆盖正在运行的内核；若设备中断号描述错误，驱动可能永远收不到 completion；若 NUMA 拓扑错误，调度和页分配会做出错误局部性判断。

\subsection{把组成原理映射到 OS 章节}

从这一节开始，后面的每个 OS 机制都应能反查到硬件根。

\begin{longtable}{p{0.2\linewidth}p{0.36\linewidth}p{0.32\linewidth}}
\toprule
硬件事实 & 推出的 OS 机制 & 后续章节 \\
\midrule
CPU 有 RIP、寄存器和栈 & 上下文保存、线程切换 & 调度与并发 \\
CPU 有特权级 & 系统调用、内核态、用户态 & 特权边界 \\
定时器可异步打断 CPU & 抢占式调度、超时 & 中断与调度 \\
MMU 按页检查权限 & 地址空间、缺页、COW & 虚拟内存 \\
cache/TLB 有维护成本 & per-CPU、TLB shootdown、屏障 & 调度、内存、锁 \\
设备通过 MMIO 暴露寄存器 & 驱动、设备文件、控制面 & 设备与 I/O \\
设备可 DMA 访问内存 & 描述符环、IOMMU、completion & 块层和网络 \\
存储设备有缓存和 flush & \code{fsync}、日志、崩溃恢复 & 文件系统和持久化 \\
\bottomrule
\end{longtable}

因此，本册合并讲组成原理和 OS 时，不追求把所有芯片细节都展开成硬件设计课，而是抓住 OS 必须依赖的组成原理主干：计算如何执行，状态如何保存，地址如何翻译，设备如何通信，中断如何进入内核，保护如何由硬件强制执行。

\subsection{最小验收：概念必须能落到机制}

读完本节后，下面这些问题应该能用自己的话讲清楚。

\begin{rubricbox}
\begin{itemize}
\tightlist
\item CPU 为什么需要程序计数器、寄存器、状态寄存器和内存接口。
\item ISA 和微架构有什么区别，为什么 OS 同时关心二者。
\item load/store 指令如何经过地址翻译、权限检查、cache 和异常路径。
\item 特权级到底限制了什么，为什么系统调用不是普通函数调用。
\item 页表为什么不只是 map，而是隔离、懒加载、共享和权限的共同根。
\item MMIO 为什么不能当普通内存，DMA 为什么需要生命周期、屏障和 IOMMU。
\item 定时器、中断控制器和精确异常如何让内核收回控制权并恢复执行。
\item cache、TLB、NUMA 和多核一致性为什么会影响调度、锁和页表修改。
\end{itemize}
\end{rubricbox}

若这些问题仍然说不清，后面看到进程、虚拟内存、驱动和文件系统时，应回到本节重新把硬件路径画出来。OS 的每个抽象都不是凭空定义，而是从硬件限制、失败场景和可验证机制中长出来的。
